\documentclass[10pt,a4paper]{amsart}
\usepackage{amsmath,amssymb,amsthm,graphicx}
% Retained minimal preamble for the Stacks Project chapter varieties.tex.
% Verbatim excerpt of the chapter preamble: Xy-pic setup (native preamble lines 14-18), the native theorem
% declarations with the mathematical macros (native preamble lines 154-217), and the project own no-op
% definitions of the reference/slogan/history display environments (native wrapper.tex lines 1-3) that the
% counted statement uses for its slogan. Copied mechanically; nothing added, removed or redefined.
% For commutative diagrams we use Xy-pic
\usepackage[all]{xy}

% We use 2cell for 2-commutative diagrams.
\xyoption{2cell}
% Theorem environments.
%
\theoremstyle{plain}
\newtheorem{theorem}[subsection]{Theorem}
\newtheorem{proposition}[subsection]{Proposition}
\newtheorem{lemma}[subsection]{Lemma}

\theoremstyle{definition}
\newtheorem{definition}[subsection]{Definition}
\newtheorem{example}[subsection]{Example}
\newtheorem{exercise}[subsection]{Exercise}
\newtheorem{situation}[subsection]{Situation}

\theoremstyle{remark}
\newtheorem{remark}[subsection]{Remark}
\newtheorem{remarks}[subsection]{Remarks}

\numberwithin{equation}{subsection}

% Macros
%
\def\lim{\mathop{\mathrm{lim}}\nolimits}
\def\colim{\mathop{\mathrm{colim}}\nolimits}
\def\Spec{\mathop{\mathrm{Spec}}}
\def\Hom{\mathop{\mathrm{Hom}}\nolimits}
\def\Ext{\mathop{\mathrm{Ext}}\nolimits}
\def\SheafHom{\mathop{\mathcal{H}\!\mathit{om}}\nolimits}
\def\SheafExt{\mathop{\mathcal{E}\!\mathit{xt}}\nolimits}
\def\Sch{\mathit{Sch}}
\def\Mor{\mathop{\mathrm{Mor}}\nolimits}
\def\Ob{\mathop{\mathrm{Ob}}\nolimits}
\def\Sh{\mathop{\mathit{Sh}}\nolimits}
\def\NL{\mathop{N\!L}\nolimits}
\def\CH{\mathop{\mathrm{CH}}\nolimits}
\def\proetale{{pro\text{-}\acute{e}tale}}
\def\etale{{\acute{e}tale}}
\def\QCoh{\mathit{QCoh}}
\def\Ker{\mathop{\mathrm{Ker}}}
\def\Im{\mathop{\mathrm{Im}}}
\def\Coker{\mathop{\mathrm{Coker}}}
\def\Coim{\mathop{\mathrm{Coim}}}

% Boxtimes
%
\DeclareMathSymbol{\boxtimes}{\mathbin}{AMSa}{"02}

%
% Macros for moduli stacks/spaces
%
\def\QCohstack{\mathcal{QC}\!\mathit{oh}}
\def\Cohstack{\mathcal{C}\!\mathit{oh}}
\def\Spacesstack{\mathcal{S}\!\mathit{paces}}
\def\Quotfunctor{\mathrm{Quot}}
\def\Hilbfunctor{\mathrm{Hilb}}
\def\Curvesstack{\mathcal{C}\!\mathit{urves}}
\def\Polarizedstack{\mathcal{P}\!\mathit{olarized}}
\def\Complexesstack{\mathcal{C}\!\mathit{omplexes}}
% \Pic is the operator that assigns to X its picard group, usage \Pic(X)
% \Picardstack_{X/B} denotes the Picard stack of X over B
% \Picardfunctor_{X/B} denotes the Picard functor of X over B
\def\Pic{\mathop{\mathrm{Pic}}\nolimits}
\def\Picardstack{\mathcal{P}\!\mathit{ic}}
\def\Picardfunctor{\mathrm{Pic}}
\def\Deformationcategory{\mathcal{D}\!\mathit{ef}}
\newenvironment{reference}{}{}
\newenvironment{slogan}{}{}
\newenvironment{history}{}{}
\usepackage{cleveref}
\crefname{originalreference}{original reference}{original references}
\Crefname{originalreference}{Original reference}{Original references}
\makeatletter
\newlabel{morphisms-definition-degree}{{02NY}{source}{}{sourcepointer-1}{}}
\newlabel{morphisms-definition-degree@cref}{{[originalreference][0][]02NY}{[1][0][]source}{}{}{}}
\makeatother
\crefname{originalreference}{original reference}{original references}
\Crefname{originalreference}{Original reference}{Original references}
\makeatletter
\newlabel{coherent-lemma-relative-affine-cohomology}{{089W}{source}{}{sourcepointer-2}{}}
\newlabel{coherent-lemma-relative-affine-cohomology@cref}{{[originalreference][0][]089W}{[1][0][]source}{}{}{}}
\makeatother
\crefname{originalreference}{original reference}{original references}
\Crefname{originalreference}{Original reference}{Original references}
\makeatletter
\newlabel{cohomology-lemma-projection-formula}{{01E8}{source}{}{sourcepointer-3}{}}
\newlabel{cohomology-lemma-projection-formula@cref}{{[originalreference][0][]01E8}{[1][0][]source}{}{}{}}
\makeatother
\crefname{originallemma}{original lemma}{original lemmas}
\Crefname{originallemma}{Original lemma}{Original lemmas}
\makeatletter
\newlabel{lemma-degree-on-proper-curve}{{0AYV}{source}{}{sourcepointer-4}{}}
\newlabel{lemma-degree-on-proper-curve@cref}{{[originallemma][0][]0AYV}{[1][0][]source}{}{}{}}
\makeatother
\title{The degree of a pulled-back locally free module along a nonconstant morphism of proper curves is multiplied by the degree of the morphism}
\author{The Stacks Project Authors}
\date{stacks-project:varieties.tex:local-snapshot-20261011; GNU FDL 1.2 or later}
\begin{document}
\maketitle
\noindent\textit{Editorial note.} All delivered mathematical text is the author's own, copied verbatim from the pinned Stacks Project chapter file varieties.tex; no statement or proof step was rewritten, reordered or removed.\par\medskip
\noindent\textit{Editorial note.} The counted span cites four cross-chapter labels (including a definition), each recorded in assembly.external\_references with its canonical tag as a navigation pointer.\par\medskip
\noindent\textit{Editorial note.} The proof cites no bibliography item inside the counted span.\par\medskip
\noindent\textit{Allowed standard prerequisites (editorial).}\begin{itemize}
\item Definition of the degree of a dominant morphism of proper curves. Stacks Project, Morphisms, Definition morphisms-definition-degree (tag 02NY). Supplies the meaning of deg(X/Y) used throughout the computation.
\item Relative affine cohomology, the projection formula and the rank formula for the degree. Stacks Project, Cohomology of Schemes, Lemma coherent-lemma-relative-affine-cohomology (tag 089W); Cohomology, Lemma cohomology-lemma-projection-formula (tag 01E8); Varieties, Lemma lemma-degree-on-proper-curve (tag 0AYV). Provide the three equalities of the displayed chain.
\end{itemize}\medskip
\ifdefined\hypertarget\hypertarget{sourcepointer-1}{}\fi
\noindent\textit{Source reference.} Original reference 02NY: sources/stacks-chow/tags.txt, source lines 3497--3497. Cross-chapter standard prerequisite cited by the author inside the counted span; it is used as a pointer and its own proof is not part of this item.\par\medskip
\ifdefined\hypertarget\hypertarget{sourcepointer-2}{}\fi
\noindent\textit{Source reference.} Original reference 089W: sources/stacks-chow/tags.txt, source lines 10511--10511. Cross-chapter standard prerequisite cited by the author inside the counted span; it is used as a pointer and its own proof is not part of this item.\par\medskip
\ifdefined\hypertarget\hypertarget{sourcepointer-3}{}\fi
\noindent\textit{Source reference.} Original reference 01E8: sources/stacks-chow/tags.txt, source lines 1825--1825. Cross-chapter standard prerequisite cited by the author inside the counted span; it is used as a pointer and its own proof is not part of this item.\par\medskip
\ifdefined\hypertarget\hypertarget{sourcepointer-4}{}\fi
\noindent\textit{Source reference.} Original lemma 0AYV: sources/stacks-chow/tags.txt, source lines 13803--13803. Cross-chapter standard prerequisite cited by the author inside the counted span; it is used as a pointer and its own proof is not part of this item.\par\medskip
\begin{lemma}
\label{lemma-degree-pullback-map-proper-curves}
Let $k$ be a field. Let $f : X \to Y$ be a nonconstant morphism of
proper curves over $k$. Let $\mathcal{E}$ be a locally free
$\mathcal{O}_Y$-module. Then
$$
\deg(f^*\mathcal{E}) = \deg(X/Y) \deg(\mathcal{E})
$$
\end{lemma}

\begin{proof}
The degree of $X$ over $Y$ is defined in
Morphisms, Definition \ref{morphisms-definition-degree}.
Thus $f_*\mathcal{O}_X$ is a coherent $\mathcal{O}_Y$-module
of rank $\deg(X/Y)$, i.e.,
$\deg(X/Y) = \dim_{\kappa(\xi)} (f_*\mathcal{O}_X)_\xi$ where $\xi$
is the generic point of $Y$. Thus we obtain
\begin{align*}
\chi(X, f^*\mathcal{E})
& =
\chi(Y, f_*f^*\mathcal{E}) \\
& =
\chi(Y, \mathcal{E} \otimes f_*\mathcal{O}_X) \\
& =
\deg(X/Y) \deg(\mathcal{E}) + n \chi(Y, f_*\mathcal{O}_X) \\
& =
\deg(X/Y) \deg(\mathcal{E}) + n \chi(X, \mathcal{O}_X)
\end{align*}
as desired. The first equality as $f$ is finite, see
Cohomology of Schemes, Lemma \ref{coherent-lemma-relative-affine-cohomology}.
The second equality by projection formula, see
Cohomology, Lemma \ref{cohomology-lemma-projection-formula}.
The third equality by Lemma \ref{lemma-degree-on-proper-curve}.
\end{proof}

\end{document}
